% CS 260D final report — skeleton (stream D owns the file; each section has an owner).
% Structure from README §10, owners from README §8.
%
% Rules:
%   - Every figure comes from analysis/plots.py or analysis/bound_check.py and
%     is \includegraphics'd from report/figures/. No hand-edited figures.
%   - Every number in the text is copied from a results-derived table.
%   - Draft complete Nov 30; one full read by someone other than the author.
%
% TODO(D): check the syllabus for a required template and page limit and
%          switch \documentclass accordingly.

\documentclass[11pt]{article}
\usepackage[margin=1in]{geometry}
\usepackage{amsmath, amssymb}
\usepackage{graphicx}
\usepackage{booktabs}
\usepackage{hyperref}

\title{HOGWILD! on Dense, Non-Convex Problems}
\author{TODO: four names}
\date{CS 260D, Fall 2026}

\begin{document}
\maketitle

\begin{abstract}
% TODO(D): write last. One sentence each: gap in the theory, what we measured,
% the three findings (H1, H2, H3), the theory check.
\end{abstract}

\section{Introduction and related work}  % owner: D
% - HOGWILD! (Niu et al. 2011): lock-free SGD, guarantees under sparsity.
% - Perturbed-iterate analysis (Mania et al. 2015).
% - Backup workers (Chen et al. 2016).
% - The lecture's open problems; which two we address (dense, non-convex) and
%   the one we touch (reproducibility).
% - Research questions H1–H3 (README §1).

\section{Setup}  % owner: B
% - Tasks T1–T3b: dataset, model, parameter count, touched-coordinate fraction.
% - Target loss per task and how it is fixed from the serial run.
% - Step-size tuning on serial.
% - Hardware.

\section{Method}  % owner: Richard
% - Shared-memory runner: process pool, pinning, one thread per worker.
% - Modes: serial, sync, backup, hogwild, hogwild with injected staleness.
% - Instrumentation: version counter and tau, injection mechanism, straggler
%   model, interference probe, evaluator snapshots.

\section{Results}

\subsection{E1: scaling}  % owner: B
% Primary figure: time-to-target vs. cores, one line per mode, one panel per task.
% Then epochs-to-target vs. m and the tau histograms.

\subsection{E2: injected staleness}  % owner: D

\subsection{E3: stragglers}  % owner: C

\subsection{Theory check}  % owner: D
% Measured slowdown vs. the perturbed-iterate prediction; state plainly where
% the theorem's assumptions do not hold.

\section{Threats to validity}  % owner: C
% - Timing resolution (2 s evaluator period), process overhead, memory bandwidth.
% - Pure-compute ceiling for speed-up (README §9).
% - Non-atomic lock-free adds; version-counter accuracy.
% - Three seeds.

\section{Conclusion}  % owner: D

\section*{Contribution statement}  % all — required by the syllabus

\appendix
\section{Implementation details}  % owner: Richard
% Start method, pinning and core map, snapshot non-atomicity, the version
% counter design, how staleness is injected, straggler calibration.

% TODO(D): references — Niu et al. 2011; Mania et al. 2015; Chen et al. 2016;
% CS 260D Lecture 4.

\end{document}
